% ==============================================================================
% secoes/03_metodologia.tex - SEÇÃO III: METODOLOGIA
% Tema: Federico Faggin e o Primeiro Microprocessador Monolítico (Intel 4004)
% ==============================================================================

\section{Metodologia}
\label{sec:metodologia}

\subsection{Abordagem Metodológica e Tipologia de Pesquisa}
\label{subsec:metodologia_abordagem}

% [DIRETIVA DE ESCRITA: TIPOLOGIA DA PESQUISA]
% - Definir a abordagem como pesquisa histórica, qualitativa e analítico-documental,
%   focalizada na interface entre a física dos semicondutores e a arquitetura de computadores.
% - Justificar o delineamento histórico-pedagógico adotado como instrumento para investigar
%   a gênese de um marco tecnológico da engenharia e suas repercussões na educação universitária.

\subsection{Coleta de Dados e Corpus Documental}
\label{subsec:metodologia_coleta}

% [DIRETIVA DE ESCRITA: CORPUS DE FONTES PRIMÁRIAS E SECUNDÁRIAS]
% - Detalhar as fontes primárias consultadas:
%   1. Patentes seminais depositadas junto ao United States Patent and Trademark Office (USPTO):
%      - Patente US 3.475.234 de Federico Faggin (depositada em 1968, concedida em 1969)
%        sobre o processo de fabricação MOS com porta de silício autoalinhada (SGT) \cite{faggin1969patent};
%      - Patente US 3.821.715 de Federico Faggin, Marcian E. Hoff e Stanley Mazor
%        (Memory System for Multi-Chip Digital Computer, concedida em 1974) \cite{faggin1974patent};
%   2. Manuais técnicos e folhas de dados originais publicados pela Intel Corporation:
%      - MCS-4 Micro Computer Set Users Manual (ed. 1971 e 1973) \cite{intel1973mcs4};
%   3. Relatórios técnicos e memórias autobiográficas de Federico Faggin
%      (notadamente "Silicon: From the Invention of the Microprocessor to the New Science of Consciousness", 2021) \cite{faggin2021}.
% - Detalhar as fontes secundárias:
%   1. Periódicos especializados em história da tecnologia e microeletrônica
%      (IEEE Annals of the History of Computing, IEEE Micro, Solid-State Electronics) \cite{faggin1970, faggin1996, cavin2012};
%   2. Tratados clássicos de circuitos integrados e sistemas VLSI (Weste & Harris, Mead & Conway) \cite{weste2010, mead1980}.
% - Detalhar as fontes pedagógicas:
%   1. Anais de conferências dedicadas à educação em engenharia (IEEE Frontiers in Education - FIE,
%      Congresso Brasileiro de Educação em Engenharia - COBENGE);
%   2. Matrizes curriculares e planos pedagógicos dos cursos de Engenharia Elétrica e de Computação da UFMG.

\subsection{Técnica de Triangulação Historiográfica e Análise Curricular}
\label{subsec:metodologia_triangulacao}

% [DIRETIVA DE ESCRITA: TÉCNICA DE TRIANGULAÇÃO]
% - Explicar a técnica de triangulação analítica:
%   * Vértice 1: Registros empíricos de patentes, diagramas esquemáticos e manuais de engenharia do MCS-4;
%   * Vértice 2: Historiografia e memórias técnicas dos protagonistas da revolução dos microprocessadores;
%   * Vértice 3: Análise pedagógica da evolução das diretrizes curriculares e da transição
%     do ensino analógico/discreto para sistemas computacionais integrados e co-design hardware/software.
